\section*{Reading Guide and Measurement Definitions}
\label{sec:supp-terms}
\newenvironment{guidelist}{\small\begin{list}{}{\setlength{\leftmargin}{0.9em}\setlength{\itemindent}{-0.9em}\setlength{\itemsep}{0pt}\setlength{\parsep}{0pt}\setlength{\topsep}{1pt}}}{\end{list}}
Part 1 maps the main paper's claims C1--C4 to their evidence, Part 2 its stated limits, Part 3 the studies and terms. Recording sites and interval conventions are in Sec.~\ref{supp:recording-overview}; a reusable protocol and recomputation checks in Sec.~\ref{sec:supp-protocol-card}. Results are for $\pi_{0.5}$ unless stated.

\paragraph{1. How this supplement supports the main paper.}
\textbf{C1: perturbations change the features along 16 shared directions per layer} (main Sec.~\ref{M-sec:shared}, Fig.~\ref{M-fig:token-structure}a).
\begin{guidelist}
\item \textbf{Sec.~\ref{sec:mean-controls}, Table~\ref{tab:mean-centered-controls}.} Held-out tasks and families, all 536 pairs: 36.7--84.2\% of the energy around the fitting-data mean, against 0.5--10.6\% (clean-feature) and 1.3--1.9\% (random); 16.26--75.00\% of the raw energy.
\item \textbf{Sec.~\ref{sec:current-basis-readout}.} The 404-pair subset of main Fig.~\ref{M-fig:token-structure}: 36.5--83.6\%, against 0.5--10.2\% and 1.3--1.9\%.
\item \textbf{Table~\ref{tab:mean-controls}.} Not the mean: uncentered shared directions retain 32.4--75.3\% of the raw energy, the span of the position means 13.5--30.4\%.
\item \textbf{Sec.~\ref{sec:rank-sensitivity}, Figs.~\ref{fig:supp-rank-joint},~\ref{fig:supp-rank-known}.} Same ordering for 1 to 32 directions.
\item \textbf{Secs.~\ref{sec:supp-energy-aggregation}, \ref{sec:supp-pi0-internals}, \ref{sec:supp-p08-dynamic}.} Energy-weighted tasks: 38.4--84.0\%. $\pi_0$: 39--77\%, against 1--28\% and about 1.5\%. A later observation (original LIBERO; tasks held out, conditions seen in fitting): 32.7--64.2\%, against 0.8--11.0\% and 1.4--1.8\%.
\item \textbf{Sec.~\ref{sec:information-budget}.} A reconstruction uses 160 coefficients from the response it scores.
\end{guidelist}
\textbf{C1, linear readout} (main Fig.~\ref{M-fig:token-structure}b).
\begin{guidelist}
\item \textbf{Sec.~\ref{sec:current-basis-readout}, Table~\ref{tab:current-basis-control}.} The 16 coordinates predict the command change better than no change at L18 in all families, worse at L1 and L9 in 7 of 8 cells. At L18 they match clean-feature coordinates (intervals span zero); first flow predicts best in every row.
\item \textbf{Sec.~\ref{sec:supp-readout-capacity}.} Nonlinear readouts change 102 and 76 of 216 interval categories.
\end{guidelist}
\textbf{C2: restoration at L9 reduces command error} (main Sec.~\ref{M-sec:restoration}, Figs.~\ref{M-fig:teaser}b,~\ref{M-fig:native-restoration}).
\begin{guidelist}
\item \textbf{Sec.~\ref{sec:allstep-restoration}, Tables~\ref{tab:allstep-restoration-contrasts}--\ref{tab:allstep-restoration-arms}.} 216 held-out cases: error 0.409 to 0.312 (Camera), 0.209 to 0.199 (Sensor Noise), lower in 25 of 27 and 24 of 30 tasks. Task success is higher by 1.85 and 3.33 points. The 16 directions reach 86.9\% and 70.7\% of replacing the whole L9 state, which gains 7.41 and 5.83 points (post hoc).
\item \textbf{Sec.~\ref{sec:dev-restoration-controls}.} At matched edit size, restoration beats clean-feature directions under Camera only.
\item \textbf{Sec.~\ref{sec:supp-native-ablation}.} No clean observation (6,048 planned rollouts): L18 response lower by 8.8\% and 7.6\%; random directions at most 0.3\%.
\end{guidelist}
\textbf{C3: command changes reach the robot} (main Sec.~\ref{M-sec:early}).
\begin{guidelist}
\item \textbf{Sec.~\ref{sec:early-late-association}, Fig.~\ref{fig:supp-early-late}.} Rank correlation with end-effector separation 0.95/0.97 after 5 controls, 0.47/0.41 after 60.
\item \textbf{Sec.~\ref{sec:supp-q0-consistency}.} Agreement across noise draws 0.88--0.92 and 0.80--0.89; lower for the gripper.
\end{guidelist}
\textbf{C4: the size of a change is not its cost} (main Sec.~\ref{M-sec:sensitivity}, Fig.~\ref{M-fig:behavior}).
\begin{guidelist}
\item \textbf{Sec.~\ref{sec:supp-matched-motion}, Table~\ref{tab:sensor-case-reference}.} $\pi_0$, Sensor Noise: 0.58 times a clean rerun's motion, 14.3 points lost; every lost case listed.
\item \textbf{Sec.~\ref{sec:five-arm}.} Lost and gained successes against the clean rerun; bounds for 39 of 2,680 missing outcomes. Lost successes exceed the rerun level under Camera and Sensor Noise, gained ones do not. AUROC 0.72/0.64 pooled, 0.62/0.55 within a family (post hoc).
\item \textbf{Secs.~\ref{supp:legacy-d3},~\ref{supp:legacy-d2}.} Rank correlation $-0.90$ across ten condition means, $-0.075$ across 400 cells. Headlight: 5.1 times the command change of matched dimming, both 77 of 80 (Fig.~\ref{fig:legacy-d2-inputs}).
\item \textbf{Sec.~\ref{sec:supp-complete-panels}.} Complete main-paper panels.
\end{guidelist}
\textbf{Additional studies, not relied on by the main paper:} internal associations (Sec.~\ref{sec:supp-bounded-internal}, Table~\ref{tab:supp-internal-links}), gripper timing (Sec.~\ref{sec:supp-landmark-final}), LIBERO-Plus outcomes (Sec.~\ref{sec:supp-plus}).

\paragraph{2. Limits stated in the main paper and where they are documented} (main Sec.~\ref{M-sec:discussion}).
\begin{guidelist}
\item \textbf{Sensitivity is measured at the episode start:} Sec.~\ref{supp:recording-overview}; the command--motion link weakens by 60 controls (Sec.~\ref{sec:early-late-association}); one later observation (Sec.~\ref{sec:supp-p08-dynamic}).
\item \textbf{Restoration: one layer, 16 directions, two families, 216 cases, needs a clean observation:} Sec.~\ref{sec:allstep-restoration}; the test without one: Sec.~\ref{sec:supp-native-ablation}.
\item \textbf{The shared directions carry the Background response, which costs $\pi_{0.5}$ no net success:} Table~\ref{tab:mean-centered-controls}; Sec.~\ref{sec:five-arm}.
\item \textbf{Simulation, two related policies, moderate settings:} Table~\ref{tab:cohorts}; Sec.~\ref{sec:supp-plus}.
\end{guidelist}

\paragraph{3. Study index.}
\begin{guidelist}
\item \textbf{Internal recordings} ($\pi_{0.5}$, LIBERO-Plus): 536 pairs on 40 tasks fit the shared directions (Secs.~\ref{sec:mean-controls}, \ref{sec:supp-plus}); 404 pairs on 28 common tasks link features, commands, motion and outcomes (Sec.~\ref{sec:supp-bounded-internal}).
\item \textbf{Restoration} ($\pi_{0.5}$): 31 held-out tasks, 108 settings, two draws, six runs, 1,296 rollouts (Sec.~\ref{sec:allstep-restoration}).
\item \textbf{Paired executions} ($\pi_{0.5}$, $\pi_0$): five runs per setting; 268 settings per policy on 40 tasks, 202 on 28 for outcomes (Sec.~\ref{sec:five-arm}). Same settings as the recordings; $\pi_0$ internals regenerated on 369 of 404 pairs (Sec.~\ref{sec:supp-pi0-internals}).
\item \textbf{Original LIBERO} ($\pi_{0.5}$): later observation, 3,360 episodes (Sec.~\ref{sec:supp-p08-dynamic}; its L9 directions are the ones removed in Sec.~\ref{sec:supp-native-ablation}); removal without a clean observation, 288 cases $\times$ 3 conditions (clean, Camera $+15^\circ$, headlight off) $\times$ 7 runs; matched lighting, 80 initial states (Sec.~\ref{supp:legacy-d2}); ten conditions, 2,344 of 2,400 pairs on 40 tasks.
\end{guidelist}

\paragraph{Glossary.}
\begin{guidelist}
\item \textbf{Setting, pair, case.} Setting: one task under one visual variant. Pair: one setting under one policy-noise draw (label 0 or 1); its clean and perturbed runs share task, start state and noise. Case: a pair in the restoration study. Start state: task, initial state and, where used, noise label.
\item \textbf{Runs.} Clean $R_0,R_1$; perturbed $T_0,T_1$ (noise of $R_0,R_1$); $R'_0$ repeats $R_0$; clean rerun: $R_1$ vs.\ $R_0$. Restoration: $N_C$ clean, $N_T$ perturbed unedited, $R_{\rm resp}$ restoration, $R_{\rm clean}$, $R_{\rm rand}$ controls, $R_{\rm full}$ whole L9 state.
\item \textbf{Directions.} Sixteen channel directions fitted to responses (shared; ``Response'' in tables), to clean features (``Clean'') or random (``Random''); position-centered fits subtract the ten position means of the fitting data.
\item \textbf{Energy shares.} One minus reconstruction error over the energy around the fitting-data mean (main Eq.~\ref{M-eq:p08-reconstruction}), or over the raw energy.
\item \textbf{Computation.} Query: one observation and plan; generation step $j=0,\ldots,9$; plan position: one planned action; control: one executed action. First five commands: their translation part, 15 numbers. First flow: clean--perturbed difference of the step-0 update.
\item \textbf{Internal quantities.} Prefix H, normalized Z, cached values V; action-expert H after a block; attention probability P; attention output W before gating (Table~\ref{tab:recording-roles}). Conditional TV: redistribution within a region. Participation: effective number of channels (tokens) carrying a response. Raw output response: a difference RMS; relative response divides it by the matched-clean norm.
\item \textbf{Command change, footprint.} Perturbed minus clean first five commands; an edit's footprint is its edited-minus-unedited change as a fraction of it.
\item \textbf{Outcomes.} Clean then perturbed (S success, F failure): SF lost, FS gained success. Success change (RD in some tables): perturbed minus clean success rate. Command--cost correlation: across the ten conditions.
\item \textbf{Intervals} resample whole tasks. \textbf{Bounds over missing outcomes} cover all feasible values of unknown outcomes.
\item \textbf{Difficulty level:} LIBERO-Plus level 1--5, higher is harder. \textbf{Camera:} Camera Viewpoints family; on original LIBERO, $+15^\circ$ unless stated.
\item \textbf{Abbreviations.} RMS root mean square; MSE mean squared error; EEF end effector; TV total variation; AUROC (AUC in some tables) area under the ROC curve; pp percentage points.
\end{guidelist}
